% Generated by report_gen.py (ai-ee v1) - do not hand-edit.
\documentclass[11pt,a4paper]{article}
\usepackage[margin=2.2cm]{geometry}
\usepackage{graphicx}
\usepackage{booktabs}
\usepackage{longtable}
\usepackage{pdfpages}
\makeatletter
\@ifundefined{KV@Gin@artifact}{\define@key{Gin}{artifact}[]{}}{}
\makeatother
\usepackage{xcolor}
\usepackage[hidelinks]{hyperref}
\setcounter{tocdepth}{1}
\begin{document}
\sloppy
\begin{center}
{\LARGE\bfseries lumina-par --- Design Document}\\[6pt]
{\large ai-ee v1 pipeline}\\[2pt]
generated 2026-07-28 19:48:21
\end{center}
\tableofcontents

\section{Board and Run Metadata}
{\small
\begin{longtable}{lp{11cm}}
\toprule
\textbf{Field} & \textbf{Value} \\
\midrule
\endhead
board & lumina-par \\
workspace & \texttt{boards/lumina-par} \\
phase & P2 \\
gate status & no gates recorded yet \\
state created & 2026-07-28T15:00:15 \\
state updated & 2026-07-28T19:47:13 \\
\bottomrule
\end{longtable}
}
\section{Overview}
\section{Requirements}
\subsection*{Requirements: lumina-par (LUM-DTR-PAR-A / LUM-PAR-A)}
Sources, in precedence order (later overrides earlier): 1. \texttt{brief/00-lumina-system-context.md} - shared system context. 2. \texttt{brief/03-rgbw-par-daughter-brief.md} - this board's brief. 3. \texttt{brief/01-carrier-board-brief.md} - the mating carrier, CONTEXT ONLY. 4. \texttt{brief/05-lumina-closed-decisions.md} - BINDING; supersedes the open-decision table in \texttt{00} s6 and corrects figures in \texttt{00} and \texttt{03}. 5. \texttt{brief/06-connector-icd.md} - FROZEN interface control document (ICD-01), owned by LUM-CAR-A. Hard input. Never redefined here.

Unstated low-risk items are marked \texttt{ASSUMED:} inline. Section 9 holds every unknown that changes the design or touches safety. Section 8 flags one \textgreater{}30 V condition that is unconditional (48 V is on the connector whether or not this board taps it) and one conditional \textgreater{}3 A condition.

Two figures in the briefs are dead and are NOT used anywhere below: the "\textasciitilde{}8.5 W sustained" daughter budget in \texttt{00} s5.1 and \texttt{03} PAR-REQ-11, and the Si3402-B / AN956 "\textasciitilde{}10 W regulated" it derives from. The carrier uses a TPS2378-class PD controller plus a 100 V buck; the re-derived budget is \textbf{8.6-9.3 W (af) / 18.7-20.0 W (at)}. PAR-REQ-11 itself still holds hard.

---

\subsubsection*{1. Function}
The colour and mood workhorse of the LUMINA fixture set: an RGBW par daughter board that stacks above the universal LUM-CAR-A carrier on the frozen ICD-01 expansion connector and turns four (possibly five) carrier PWM channels into four independent constant-current LED drives. It runs continuously rather than in bursts, so its problem is not peak power but quality of dimming and consistency between the 6-8 fixtures in the room: visually stepless dimming at low output, slow hue drift over tens of seconds with no banding, and no per-fixture tint variation in a synchronised full-room wash. It carries LED drivers, thermal sensing, an independent over-temperature shutdown, the ENABLE gate, the board ID leg and the calibration store - and nothing else. It has no MCU, no network interface, no external connector of any kind, and no energy store for burst output (that is the strobe board's problem, not this one's).

---

\subsubsection*{2. Interfaces}
\textbf{Every interface on this board is through ICD-01. There are no others.} ICD-01 s9 is explicit: a daughter may not add an external connector of any kind - no barrel jack, no DMX, no second Ethernet. Ethernet on the carrier is the fixture's only external connection, and that is what makes the non-isolated PoE topology compliant.

\paragraph*{2.1 J3 - POWER block, 2x7, daughter side}
\begin{itemize}
\item Part: {*}\emph{CONNFLY DS1023-2}7SF11{*}{*} socket (14 pos, 2.54 mm, 600 V,
3 A/contact, 8.5 mm body). Second sources on the same land pattern:
HanElectricity 2541FV, Boomele 2.54-2{*}7P, HCTL HC-PZ254. JLC Extended -
there is no JLC Basic part in any board-to-board family, so "prefer Basic"
is unachievable here and is not a selection error.
\item \textbf{Reverse-mounted on the daughter's BOTTOM side, facing down} (ICD s7.3).
\item Pin map (ICD s3.1), authoritative, do not re-derive:
1/3/5 \texttt{+48V\_SW}; 2/4/6/7/8/10/13 \texttt{GND}; 9/11 \texttt{+12V}; 12/14 \texttt{+3V3}.
Position 1 silkscreen-marked with a triangle.
\item Rail order along the connector is 48 -\textgreater{} GND -\textgreater{} 12 -\textgreater{} 3.3 so that no
single-position mis-seat can put a higher rail on a lower rail's pin.
\end{itemize}
\paragraph*{2.2 J4 - SIGNAL block, 2x12, daughter side}
\begin{itemize}
\item Part: {*}\emph{CONNFLY DS1023-2}12SF11{*}{*} socket (24 pos, 600 V, 3 A/contact,
8.5 mm body), same mounting rules and same second-source situation.
\item Pin map (ICD s3.2), authoritative: \texttt{PWM0..7} at positions 1,2,5,6,7,8,11,12;
\texttt{GND} at 3,4,9,10,13; \texttt{DSPI\_SCK} 14, \texttt{MOSI} 15, \texttt{MISO} 16, \texttt{CSn} 17;
\texttt{I2C\_SCL} 18, \texttt{SDA} 19; \texttt{ADC0} 20, \texttt{ADC1} 21; \texttt{ID\_ADC} 22; \texttt{ENABLE} 23;
\texttt{FAULT} 24. \textbf{No 48 V exists anywhere on this connector.}
\end{itemize}
\paragraph*{2.3 Signal contract (ICD s3.3) - what this board must honour}
\begin{quote}\small\ttfamily\raggedright
\textbar{} Signal \textbar{} Requirement on this board \textbar{}\\
\textbar{}---\textbar{}---\textbar{}\\
\textbar{} `PWM0..7` \textbar{} 3.3 V push-pull, {*}{*}13-bit at 9.766 kHz default{*}{*} (CAR-REQ-12). PWM0-3 = LEDC timer 0, PWM4-7 = timer 1; channels on a timer share frequency AND resolution. `ASSUMED:` R,G,B,W map to PWM0-3 (timer 0); a 5th channel (question 4) takes PWM4 on timer 1, programmed identically. LEDC timers 2 and 3 are unallocated and available on request - that is the only sanctioned route to a different PWM frequency for this daughter, and it needs the carrier owner's agreement, not a unilateral change. \textbar{}\\
\textbar{} `I2C\_SCL/SDA` \textbar{} 400 kHz open drain. {*}{*}Pull-ups are on the carrier (4.7 k). This board must NOT fit its own.{*}{*} The daughter owns the whole address space; the carrier reserves nothing - so EEPROM and any digital temp sensor must be allocated non-colliding addresses here. \textbar{}\\
\textbar{} `ADC0`, `ADC1` \textbar{} daughter -\textgreater{} carrier, 0-3.3 V, {*}{*}source impedance \textless{}= 10 kohm{*}{*}. This is the PAR-REQ-12 thermistor path; an NTC divider must be sized to meet the 10 kohm ceiling at every temperature of interest. \textbar{}\\
\textbar{} `ID\_ADC` \textbar{} The carrier fits the TOP leg (10 k to +3V3); {*}{*}this board fits the BOTTOM leg to GND{*}{*}. The code value is allocated by the carrier owner, not chosen here - see question 15. \textbar{}\\
\textbar{} `ENABLE` \textbar{} carrier -\textgreater{} daughter, push-pull, {*}{*}active HIGH{*}{*}. This board must fit its own {*}{*}100 kohm pull-down{*}{*}, gate {*}{*}every{*}{*} output stage with it (driver EN pins, any gate driver, any charge path), and {*}{*}never latch it locally{*}{*}. A carrier PWM pin can glitch \textasciitilde{}60 us at power-up; ENABLE is what makes that a no-op. \textbar{}\\
\textbar{} `FAULT` \textbar{} daughter -\textgreater{} carrier, {*}{*}open drain, active low{*}{*}, 10 k pull-up on the carrier, wire-OR'd with the carrier's eFuse fault. {*}{*}Never drive it high.{*}{*} \textbar{}\\
\textbar{} `DSPI\_{*}` \textbar{} Available, \textless{}= 26 MHz, mode 0, shared bus with one CS. Not expected to be needed by a 4-channel PWM design; a daughter needing more devices decodes locally. \textbar{}
\end{quote}
\paragraph*{2.4 LED interface}
Emitters are the board's real load and their mounting is \textbf{not settled} - on-board emitters versus an off-board module on a separate heatsink with an internal wiring harness is question 5, and it interacts with ICD s9's external-connector ban (an internal harness inside the enclosure is not an external connector, but this needs confirming, not assuming). Optics (PAR-REQ-14 wide wash beam from a 2.5 m ceiling, PAR-REQ-15 diffusion sufficient that R/G/B/W mix before reaching a surface) may or may not be in this run's scope - question 8.

\paragraph*{2.5 Status / test points}
No requirement stated. \texttt{ASSUMED:} no user-visible indicators on this board - CAR-REQ-09 status indication lives on the carrier. Any test point fitted must carry the ICD s9 bench-hazard warning: the entire fixture floats at PoE potential, and an earthed scope probe or non-isolated USB-UART adapter does not merely risk damage, it \textbf{breaks PD signature detection outright} (detection currents are a few hundred microamps).

---

\subsubsection*{3. Power}
\paragraph*{3.1 The envelope (binding)}
\begin{quote}\small\ttfamily\raggedright
\textbar{} \textbar{} 802.3af (first build) \textbar{} 802.3at (upgrade path) \textbar{}\\
\textbar{}---\textbar{}---\textbar{}---\textbar{}\\
\textbar{} {*}{*}Sustained power to this daughter, all rails combined{*}{*} \textbar{} {*}{*}8.6-9.3 W{*}{*} \textbar{} {*}{*}18.7-20.0 W{*}{*} \textbar{}
\end{quote}
D-01 is closed: the carrier's PD front end, converter, magnetics and thermal design are sized for \textbf{Type 2 (at)}, but classification is programmed to \textbf{Type 1 (af)} for the first build via a class resistor. The upgrade is a resistor change plus a PoE+ switch, \textbf{no respin}. D-01 says explicitly: {*}do not allow any other component to become the part that pins the design to Type 1.{*} On this board that clause lands on the LED drivers, the emitters and the heatsink - see question 2, which is the highest-consequence open item in this document.

\paragraph*{3.2 Per-rail ceilings (ICD s6.2) - individual ceilings, the TOTAL binds}
\begin{quote}\small\ttfamily\raggedright
\textbar{} Rail \textbar{} Sustained af \textbar{} Sustained at \textbar{} Peak (ms, local bulk) \textbar{} Hardware fault ceiling \textbar{}\\
\textbar{}---\textbar{}---\textbar{}---\textbar{}---\textbar{}---\textbar{}\\
\textbar{} `+48V\_SW` \textbar{} 0.25 A (12.0 W) \textbar{} 0.50 A (24.0 W) \textbar{} 1.0 A \textbar{} 1.0 A eFuse, {*}{*}latch off{*}{*} \textbar{}\\
\textbar{} `+12V` \textbar{} 0.75 A (9.0 W) \textbar{} 1.25 A (15.0 W) \textbar{} 2.0 A \textbar{} 2.0 A converter OCP \textbar{}\\
\textbar{} `+3V3` \textbar{} 0.25 A (0.83 W) \textbar{} 0.25 A \textbar{} 0.50 A \textbar{} 1.0 A converter, limit \textgreater{}= 1.3 A \textbar{}
\end{quote}
{*}{*}Derived, and load-bearing: the +12V rail alone cannot cover this board's envelope at either class.{*}{*} At af it needs 0.72-0.78 A against a 0.75 A ceiling (right at the edge); at at it needs 1.56-1.67 A against a \textbf{1.25 A} ceiling, which it misses by 3.7-5.0 W. ICD s6.3 is unambiguous: {*}anything above 1.25 A at the at operating point must be taken on \texttt{+48V\_SW}{*}, and the 1.25 A ceiling is a thermal limit on the carrier's 48-\textgreater{}12 converter in a sealed box, not a current rating that can be argued with. Taking power on \texttt{+48V\_SW} instead of \texttt{+12V} is also worth \textbf{0.67 W (af) / 1.30 W (at)} of extra delivered power because it skips a conversion.

\texttt{ASSUMED:} for an \textbf{af-sized} board, LED power is taken from \texttt{+12V}, which is D-02's stated purpose for that rail (so 6-8 fixtures do not each duplicate a \textgreater{}= 60 V converter). If question 2 answers "size for at", that assumption is void and P1 must take at least part of the load from \texttt{+48V\_SW}, inheriting the bleed path, the 0.60 mm creepage regime, 100 V capacitors, 0805-or-larger resistors across the 48 V domain, and the inrush obligation below.

\paragraph*{3.3 Hard rules inherited from the ICD}
\begin{itemize}
\item \textbf{Hardware backstop for PAR-REQ-11.} The clamp is a firmware rule, but
firmware fails. With ENABLE asserted and every PWM stuck at 100 \% (hung MCU;
the 2 s watchdog fade cannot help if the MCU is the thing that hung), the
board's total draw must remain {*}{*}below the carrier's hardware fault
ceilings{*}{*} - 12 V OCP 2.0 A, 48 V eFuse 1.0 A latch-off, PSE overload timer
\textasciitilde{}50-75 ms. A board that can only meet the budget when firmware is healthy
will brown out or latch off the whole fixture. This sizes the per-channel
drive current and is not optional.
\item \textbf{`+48V\_SW` is dead at power-up and stays dead for hundreds of ms.} The
carrier's load switch is a compliance part (802.3 caps PD port capacitance at
\textasciitilde{}180 uF); it closes only after firmware asserts ENABLE. Design for it.
\item \textbf{No first-mate/last-mate sequencing exists.} The board must tolerate 48 V
arriving before or after 3.3 V, in either order.
\item {*}{*}No path may energise anything on this board from \texttt{+12V} or \texttt{+3V3} while
\texttt{+48V\_SW} is off{*}{*}, or the 802.3 inrush compliance is defeated behind the
switch's back.
\item \textbf{Inrush is this board's responsibility} (CAR-REQ-14). Size the limiter
against the \textbf{PD's 1.0 A operating current limit}, not the connector's
5.4 A rating; sizing against the connector is the classic way to trip the
PD's 800 us foldback deglitch and brown out the fixture.
\item \textbf{If any 48 V net is tapped: a bleed path is mandatory} (CAR-REQ-17 /
STR-REQ-10). The carrier deliberately fits no series diode on \texttt{+48V\_SW}, so
the daughter's bleed is not stranded above the carrier's 100 k.
\item {*}{*}0.60 mm outer-layer copper clearance around every 48 V net, board-wide,
is NOT inherited automatically{*}{*} - this board's DRC must be set up for it at
P5 (IPC-2221B B2, 51-100 V band, 57 V worst case). The 0.13 mm coated column
is not claimable: LPI soldermask is not a qualified conformal coating and
\texttt{check\_creepage.py} implements only the uncoated columns, so a layout
designed to 0.13 mm fails P8 with no waiver mechanism.
\end{itemize}
\paragraph*{3.4 LED-side budget (ESTIMATES, clearly marked as such)}
GUESS, for question sharpening only, to be replaced by a real power table at the P1/P2 design gate (\texttt{03} review gate 1):

\begin{itemize}
\item Daughter housekeeping (driver quiescent, ID/EEPROM, sense): \textbf{\textasciitilde{}0.3-0.5 W}.
\item Switching CC driver efficiency: \textbf{\textasciitilde{}88-92 \%}.
\item Therefore electrical power reaching the emitters: \textbf{\textasciitilde{}7.6-8.5 W (af)} /
\textbf{\textasciitilde{}16.5-18.4 W (at)}.
\item A typical 4-in-1 RGBW multi-die emitter (R AlInGaP \textasciitilde{}2.1-2.9 V, G/B InGaN
\textasciitilde{}3.0-3.6 V, W \textasciitilde{}2.9-3.4 V) draws roughly \textbf{4.0-4.7 W at 350 mA/die} and
\textbf{8.4-9.8 W at 700 mA/die} with all four dies on. So \textbf{one} such emitter at
\textasciitilde{}600-700 mA/die consumes the entire af envelope by itself; the at envelope
buys roughly double.
\item Note for P1, not a decision here: PAR-REQ-10 forbids burning the R-vs-GBW
forward-voltage spread in a shared linear element, and a 12 V rail against a
2.1-3.6 V single die is a 3.3-5.7x step-down. That combination rules out
shared linear current control and points at either per-channel switching
regulators or multi-die series strings per channel. It also interacts with
question 3 (package choice).
\end{itemize}
No battery, no charging, no mains anywhere in this system.

---

\subsubsection*{4. Environment}
\begin{itemize}
\item Deployment: indoor, basement/garage room \textasciitilde{}5 m x 7 m x 2.5 m, 8-12 fixtures,
ceiling-mounted at \textasciitilde{}2.5 m. \texttt{ASSUMED:} ceiling mount, beam pointing into the
room; no vibration or shock requirement; no ingress rating (question 6
decides sealed vs vented, which is a thermal question, not an IP one).
\item \texttt{ASSUMED:} external ambient 0-40 C. Not stated anywhere in the briefs.
\item \textbf{Internal enclosure air reaches 56 C (af) / 69 C (at)} per ICD s7.6. This
is the ambient the emitters and drivers actually see, not room temperature.
\item \textbf{Continuous duty with no relief.} Unlike the strobe, this board runs hot
all the time (\texttt{03} review gate 4, and H1-Q5 restates it). Every thermal
number must be a steady-state number.
\item Heat this board must reject (ESTIMATE: LEDs convert \textasciitilde{}20-35 \% of electrical
input to light, the rest to heat): \textbf{\textasciitilde{}6.5-8.4 W (af)} / \textbf{\textasciitilde{}14-17 W (at)},
plus the carrier's own 2.4 W (af) / 3.7 W (at) in the same box.
\item \textbf{This is the riskiest number in the document.} With internal air at 69 C
and a red-die junction that wants to stay near 85-100 C for colour and
lifetime stability, the junction-to-internal-air thermal budget at the at
operating point is roughly \textbf{1.0-1.1 C/W}, and at af roughly 2.3 C/W. A
sealed plastic box has no convective path out. Questions 6 and 7 exist
because the at case very likely does not close without either ventilation or
an external heat path, and H1-Q5 constrains both (non-conductive enclosure,
heatsink NOT user-accessible).
\item Connector contacts are rated -40 to +105 C, so they are not the limit.
\end{itemize}
---

\subsubsection*{5. Size \& mounting}
All of the following is the \textbf{common LUMINA footprint} from ICD s7.1, which every board inherits. It is proposed, not yet confirmed with the human - MECH-02 requires the outline to be closed at H1 \textbf{before P5}, because \texttt{board\_init.py} binds the outline permanently and ai-ee has no outline-shrink step. See question 12.

\begin{quote}\small\ttfamily\raggedright
\textbar{} Item \textbar{} Value \textbar{}\\
\textbar{}---\textbar{}---\textbar{}\\
\textbar{} Outline \textbar{} {*}{*}100.0 x 80.0 mm{*}{*} \textbar{}\\
\textbar{} Corner radius \textbar{} {*}{*}3.0 mm{*}{*}, all four corners \textbar{}\\
\textbar{} Mounting holes \textbar{} {*}{*}4x M3 (3.2 mm) at 5 mm inset{*}{*} (a 90 x 70 mm rectangle) {*}{*}plus a 5th M3 at (46, 74){*}{*} \textbar{}\\
\textbar{} Thickness \textbar{} 1.6 mm \textbar{}\\
\textbar{} Coordinate origin \textbar{} board top-left, x right, y down \textbar{}\\
\textbar{} Mated board-to-board height \textbar{} {*}{*}11.0 mm{*}{*} hard-seated \textbar{}\\
\textbar{} Standoffs \textbar{} 5x M3 female-female, 11.0 mm \textbar{}
\end{quote}
\paragraph*{5.1 The RJ45 notch - hard requirement from P2 onward}
\textbf{A 30 x 26 mm notch in the TOP edge, region (6, 0) - (36, 26).} The carrier's board-edge magjack is \textasciitilde{}15 mm tall against an 11.0 mm stack, so it protrudes \textasciitilde{}4 mm above this board's underside. The outline rectangle, corner radius and 5-hole pattern are unchanged; only this local relief differs. It is also the \textbf{primary anti-180-degree interlock} (CAR-REQ-16): rotated, the notch lands at the bottom edge and the board presents solid material over the jack, so the boards cannot be forced flat. That is a mechanical stop, not a warning.

\paragraph*{5.2 Exclusion zones (ICD s7.6) - all four apply to this board}
\begin{quote}\small\ttfamily\raggedright
\textbar{} Zone \textbar{} Region \textbar{} Rule \textbar{}\\
\textbar{}---\textbar{}---\textbar{}---\textbar{}\\
\textbar{} RJ45 relief \textbar{} (6, 0) - (36, 26) \textbar{} cut away entirely (s5.1) \textbar{}\\
\textbar{} DC-DC hot zone \textbar{} (2, 46) - (36, 68) \textbar{} {*}{*}no LED drivers and no aluminium electrolytics{*}{*} - the carrier's 48-\textgreater{}12 converter dissipates up to 1.25 W directly below, and in a stacked mezzanine this board's parts sit vertically over it \textbar{}\\
\textbar{} Antenna column \textbar{} (88, 25) - (100, 55) \textbar{} {*}{*}no copper on any layer, no metal component{*}{*} - the carrier's ESP32-S3 PCB antenna is 11 mm below and Wi-Fi is a supported control path (H1-Q8) \textbar{}\\
\textbar{} Recovery header \textbar{} (76, 0) - (98, 20) \textbar{} keep clear enough for a 6-way jumper lead with this board fitted, or accept board removal for firmware recovery \textbar{}
\end{quote}
\paragraph*{5.3 Connector positions - PROVISIONAL, and this board is blocked on them}
ICD s7.2 gives J3 body (14, 68)-(34, 78) with position 1 at (15.3, 69.3), H5 at (46, 74), J4 body (56, 68)-(88, 78) with position 1 at (57.3, 69.3). {*}{*}ICD s7.2 is the one section not frozen at H1{*}{*}: the carrier owner must compare placed positions after its own P6, correct them, and re-issue. Daughters are explicitly blocked on that re-issue. See question 16.

Mating geometry to verify at P4/P6, not to re-derive: this board's sockets are \textbf{bottom-side, facing down}, and mate with a top-side header on the carrier. Plan-view (x, y) coordinates are shared between the two boards, but pin numbering mirrors across the row on a bottom-side footprint - pin 1 alignment must be checked in the mated view, not assumed from the footprint.

\paragraph*{5.4 Pipeline facts that constrain P5 (verified in this repo, not assumed)}
\begin{itemize}
\item \texttt{board\_init.py --mounting-holes N} places holes at \textbf{inset = margin / 2}
(confirmed in the script). The ICD's 5 mm inset therefore requires
\textbf{`--margin 10`}, not the default 6. The corner radius is separately clamped
in the worker; read the report's \texttt{corner\_radius} field and \texttt{worker\_notes}
rather than assuming the requested 3.0 mm was honoured.
\item \textbf{`board\_init.py` has NO notch/cutout option} (\texttt{--outline} takes only \texttt{auto}
or \texttt{WxH}), and \texttt{kc.py} exposes no outline-editing subcommand. The mandatory
30 x 26 mm relief must therefore be produced by a direct Edge.Cuts edit of
the \texttt{.kicad\_pcb} after \texttt{board\_init}, with every downstream consumer
(\texttt{place\_{*}}, \texttt{planes\_gen}, \texttt{dfm\_check}, \texttt{fab\_export}) re-verified against the
modified outline. \textbf{Flagged as an implementation risk for P5}, not a
requirements question - the notch itself is non-negotiable.
\item Placement keepouts ARE supported: \texttt{constraints.json} accepts
\texttt{placement.keepouts: {[}\{"rect": {[}x1,y1,x2,y2{]}, "reason": ...\}{]}} and the
annealer treats them as fixed obstacles. All four ICD s7.6 zones plus the
notch region must be declared there.
\end{itemize}
\paragraph*{5.5 Height}
Board-to-board is 11.0 mm. \textbf{Total height above this board is unknown} and depends on the emitter/heatsink/optic stack (questions 5-8). Not stated anywhere in the briefs.

---

\subsubsection*{6. Quantity \& budget}
\begin{itemize}
\item \textbf{Quantity: 6-8} par fixtures out of 8-12 total (\texttt{03} s1). Exact count and
spares not stated - question 14.
\item System budget: \textbf{\$500-1000 for fixtures + switch + enclosures} (\texttt{00} s1).
ICD s1.1 rejects a Samtec pair at \$15.30 "over half the \$30/board target",
so a \textbf{\$30/board} figure exists in the carrier run - but no unit cost target
has been stated for \emph{this} board. With 8-12 fixtures, a managed PoE+ switch
(note: D-01's upgrade path implies PoE+, which is materially more expensive
than the PoE switch the original budget assumed), enclosures, carriers,
strobes and pars all inside \$500-1000, this is tight. Question 14.
\end{itemize}
---

\subsubsection*{7. Assembly}
\begin{itemize}
\item \texttt{ASSUMED:} JLCPCB fabrication and PCBA, consistent with every other board in
this repo. Not stated in the LUMINA briefs.
\item {*}{*}Both ICD connectors are 2.54 mm THT sockets that must be reverse-mounted on
the bottom side.{*}{*} That is not a normal JLC PCBA process, so the realistic
routes are hand-soldering them post-PCBA or finding a bottom-side SMD
equivalent on the same land pattern (ICD s7.3 permits "a bottom-side SMD
equivalent"). Question 13.
\item High-power emitters, if on-board, need a large thermal pad and a via farm on
1.6 mm FR4 - or an MCPCB, which the common footprint does not provide.
Interacts with questions 5 and 13.
\item Sidedness: bottom side already carries the two sockets, so the board is
double-sided by construction. Question 13 confirms the process.
\end{itemize}
---

\subsubsection*{8. Compliance / safety flags}
\begin{quote}\small\ttfamily\raggedright
\textbar{} Flag \textbar{} Status \textbar{}\\
\textbar{}---\textbar{}---\textbar{}\\
\textbar{} Mains voltage \textbar{} {*}{*}No.{*}{*} PoE only. \textbar{}\\
\textbar{} Batteries \textbar{} {*}{*}No.{*}{*} No battery, no charging anywhere in the system. \textbar{}\\
\textbar{} Motors \textbar{} {*}{*}No.{*}{*} \textbar{}\\
\textbar{} {*}{*}\textgreater{} 30 V{*}{*} \textbar{} {*}{*}YES, unconditionally.{*}{*} `+48V\_SW` is present on J3 whether or not this board taps it. PD rail 37-57 V nominal, {*}{*}57 V worst case{*}{*} (IEEE 802.3 PSE maximum). Consequences: 0.60 mm outer-layer clearance around every 48 V net board-wide; 100 V capacitors (63 V is not enough at 57 V once ceramic DC-bias derating applies); any resistor across the 48 V domain must be 0805 or larger, or split in series; the clearance applies through the board (inner 0.10 mm / outer 0.60 mm under any 48 V antipad); mandatory bleed path if tapped. {*}{*}Insulation class is functional only{*}{*} - 57 V is below the IEC 62368-1 ES1 limit of 60 V DC, so no personnel safeguard and no safety-mandated creepage applies, and an unearthed PD needs {*}{*}no MOV-to-earth surge network{*}{*} (do not copy one from a reference design). \textbar{}\\
\textbar{} High current (\textgreater{} 3 A) \textbar{} {*}{*}Conditional - architecture must report and re-flag.{*}{*} No connector rail exceeds 1.25 A. But LED-side per-channel current is unbounded by the ICD: a single low-Vf red channel sized near the whole envelope (e.g. 8 W into a \textasciitilde{}2.4 V string) would be \textasciitilde{}3.3 A. P1/P2 must publish per-channel drive currents; if any exceeds 3 A this flag becomes unconditional. \textbar{}\\
\textbar{} RF transmit \textbar{} {*}{*}No radio on this board{*}{*} - but the carrier's ESP32-S3 radio is a {*}supported control path{*} (H1-Q8), so the (88, 25)-(100, 55) antenna column is a hard keepout (no copper on any layer, no metal component) and switching LED drivers 11 mm above a 2.4 GHz PCB antenna are a desense risk that layout sign-off must address. \textbar{}\\
\textbar{} {*}{*}Floating PoE potential{*}{*} \textbar{} {*}{*}Applies to everything this board touches.{*}{*} The entire fixture is non-isolated and floats at PoE potential: this board, its drivers, its LED wiring, and any heatsink the LED module sits on. 802.3 compliance is achieved {*}only{*} by there being no accessible external conductor. This is ICD s9 / `decisions.md` OPEN-C and it is {*}{*}unresolved{*}{*} - questions 5 and 7. \textbar{}\\
\textbar{} {*}{*}Continuous hot surface{*}{*} \textbar{} High-power LEDs at 8.6-20 W continuous inside a non-conductive enclosure whose internal air already reaches 56-69 C. Enclosure touch temperature and plastic deformation are real; H1-Q5 requires the heatsink to be non-user-accessible. \textbar{}\\
\textbar{} {*}{*}Photobiological (IEC 62471){*}{*} \textbar{} A high-power RGBW source with a wash optic is plausibly Risk Group 1-2 at close range (blue-light hazard). Not a certification requirement for a private installation, but worth a note in the design doc rather than discovering it by looking into a fixture. \textbar{}\\
\textbar{} Certification \textbar{} `ASSUMED:` private installation, own use - no CE/FCC/UL conformity assessment path required. \textbar{}
\end{quote}
\textbf{Nothing in section 8 is guessed.} The \textgreater{}30 V condition and the floating-PoE condition are stated facts from the ICD; the \textgreater{}3 A condition is explicitly left open for the architecture to answer.

---

\subsubsection*{9. Open questions}
Questions 1-4 are the four the orchestrator already queued, sharpened with the numbers above. Questions 5-16 are everything else this design genuinely needs. Each offers a recommendation. Questions \textbf{2, 6, 7 and 12} are the ones that change the board most, and 12 and 16 must be answered before P5 regardless.

\textbf{1. Total-power clamp policy (PAR-REQ-11).} Four channels cannot all run at 100 \% - the sum of four channels at full drive exceeds the 8.6-9.3 W (af) envelope by roughly 30-60 \% on any sane sizing. When the host asks for more than the fixture can deliver, what should the fixture protect? (a) \textbf{Preserve hue}: scale all four channels by the same factor. The colour is exactly right; the fixture is simply dimmer. White (all four on) is the worst case and gets dimmest. (b) \textbf{Preserve brightness}: hold the total at the cap and reallocate between channels. Brightness is stable; the hue drifts toward whichever channel is cheapest in watts, and drifts \emph{differently} per fixture, which is exactly the artifact PAR-REQ-06 exists to prevent. (c) Hybrid: preserve hue, let the white channel absorb the deficit - this desaturates at high intensity, the "visible artifact, not graceful degradation" the brief warns about. \textbf{RECOMMEND (a).} Uniform hue-preserving scaling. Colour consistency across 6-8 fixtures washing one wall is this board's reason to exist; a uniform brightness loss is invisible, a hue split across fixtures is not. Note that whichever answer is chosen, the \textbf{hardware backstop in s3.3 is still required}

\begin{itemize}
\item the board must be electrically incapable of exceeding the carrier's fault
\end{itemize}
ceilings with every PWM stuck at 100 \%.

\textbf{2. Size the LED engine for af or at?} D-01 kept the \emph{carrier} at-capable and af-classified, and told us not to let any component pin the design to Type 1. Does that clause bind this board's LED engine too? (a) \textbf{Size for af (8.6-9.3 W)}: smaller emitter, smaller heatsink, cheaper, thermally plausible in a plastic box - but a PoE+ upgrade then buys nothing on the par because the emitters and heatsink cannot use the extra \textasciitilde{}10 W, and getting it later is a respin of this board. (b) \textbf{Size for at (18.7-20.0 W), run clamped at af}: the upgrade is a firmware constant. But it roughly doubles heatsink area and emitter cost per fixture, and - derived above - {*}{*}it forces at least 3.7-5.0 W to be taken from \texttt{+48V\_SW}{*}{*} because the +12V rail tops out at 15.0 W, which means a \textgreater{}= 60 V input stage on this board after all, the thing D-02's 12 V rail was created to avoid on 6-8 fixtures. {*}{*}RECOMMEND (a) for the first build, with a written note in the design doc that the par does not follow the carrier's at upgrade.{*}{*} The par is a sustained wash, not the fixture that needs the headroom; the strobe is. Spending the at budget on the par costs money and heatsink on all 6-8 units to buy brightness the profiles ask of the strobe. But this is a lighting-design call, not an engineering one - if a 9 W par reads as underwhelming, (a) is a respin.

\textbf{3. LED package: four discrete emitters or an integrated RGBW multi-die?} (a) \textbf{Integrated 4-in-1} (single package, four dies on one thermal slug): mixes better optically and largely solves PAR-REQ-15 fringing; one thermal interface; one binning transaction. All four dies share one junction temperature, so the hottest channel heats the others, and red (AlInGaP) is both the most thermally sensitive and the one that gets cooked. (b) \textbf{Four discrete emitters}: independent thermal management, cheaper and easier sourcing, per-channel binning - but four separate optical sources need real diffusion to avoid visible R/G/B/W shadow fringing on a wall (PAR-REQ-15), and four packages spread over the board is four thermal paths to solve. \textbf{RECOMMEND (a) integrated 4-in-1}, on the grounds that PAR-REQ-15 is a hard optical requirement and mixing is the one thing an integrated package gives for free. Note the interaction with s3.4: a 12 V rail against a single 2.1-3.6 V die is a 3.3-5.7x step-down, which pushes toward per-channel switching drivers either way.

\textbf{4. A fifth channel (amber)?} PWM budget is free - the connector carries 8 channels and the par uses 4, and PWM4 sits on LEDC timer 1 which can be programmed to the same 13-bit/9.766 kHz. So the cost is not connector budget, it is: one more driver, and a fifth channel drawing from the \textbf{same} 8.6-9.3 W envelope, making the question-1 clamp bite harder in every mixed colour. P3 French melodic leans on gold/amber and RGBW amber is a mix rather than a true emitter. Caveat worth weighing: amber AlInGaP loses output at roughly -0.5 to -1 \%/C of junction temperature, roughly twice as badly as the InGaN channels, in a box whose internal air is already at 56-69 C continuously - so a calibrated amber will drift visibly as the fixture warms up unless it is temperature-compensated in firmware. \textbf{RECOMMEND: no fifth channel on rev A.} Add it only if question 2 answers "size for at" (where there is power to spare). At af it dilutes the envelope across five channels and adds the most thermally unstable emitter in the catalogue to a board that already has a thermal problem.

{*}{*}5. LED mounting and wiring: on-board emitters, or an off-board LED module on a separate heatsink?{*}{*} (a) \textbf{On-board}: emitters on this PCB's top side, heat through thermal vias into a heatsink bolted to the board. No wiring at all, so ICD s9 is trivially satisfied - but 1.6 mm FR4 is a poor thermal path, the common footprint gives no MCPCB, and the emitters' position is then dictated by the connector/keepout geometry rather than by the optics. (b) \textbf{Off-board module} on its own MCPCB/heatsink, connected by an internal wiring harness (5-9 conductors). Best thermal and optical freedom. {*}{*}ICD s9 bans a daughter from adding "an external connector of any kind" - a harness that never leaves the enclosure is not an external connector on any reasonable reading, but this needs confirming, not assuming{*}{*}, because if it is read strictly then option (b) is illegal and option (a) is the only route. {*}{*}RECOMMEND (b), with explicit confirmation that an internal wire-to-board connector (e.g. a JST-family header) is permitted{*}{*} on the grounds that ICD s9's intent is "nothing penetrates the enclosure wall", and s9 itself already anticipates the off-board case ("if the LED module is on a separate heatsink (Q4a default), that heatsink and its wiring are at PoE potential too"). If the answer is a strict no, please say so plainly - it forces (a) and changes the thermal design completely.

\textbf{6. Is the enclosure sealed, or vented?} ICD s7.6 says "sealed box" with internal air at 56 C (af) / 69 C (at); H1-Q5 says plastic, non-conductive, with the heatsink not user-accessible. No ingress requirement is stated anywhere. This matters more than any other mechanical question: from s4, a \emph{sealed} plastic box gives a junction-to-internal-air budget of \textasciitilde{}2.3 C/W at af and \textasciitilde{}1.0 C/W at at, and the at case very likely does not close. (a) Fully sealed (no openings at all). (b) \textbf{Vented} - louvres or a labyrinth sized so no finger or 4 mm probe reaches live parts or the heatsink, preserving H1-Q5's non-accessibility while allowing convection. (c) Sealed but with the heatsink forming part of the enclosure wall (see question 7). \textbf{RECOMMEND (b) vented with a finger-guard geometry.} It is an indoor basement/garage installation with no stated dust or moisture requirement, and convection is close to free. Sealing costs real money in heatsink area and caps the fixture's output.

\textbf{7. Is an external or enclosure-integrated heat path allowed?} ICD s9 and \texttt{decisions.md} OPEN-C leave this open and it is the sharp end of the non-isolated topology: any heatsink bonded to this board's ground {*}{*}floats at PoE potential{*}{*}, and "if the heatsink is touchable, metal, or shares a mount with anything earthed, the non-isolated topology is non-conformant." (a) Heatsink entirely internal, no metal reaches the outside. Thermally worst, compliance-simplest. (b) Heatsink forms part of the enclosure but is shrouded behind a plastic guard - not touchable, and not sharing a mount with anything earthed. (c) Exposed metal heatsink - \textbf{would break the compliance argument} unless the whole fixture is re-isolated, which is out of scope. \textbf{RECOMMEND (b)}, plus a hard rule that the ceiling mount is non-conductive and bonded to nothing. Also please confirm: should the LED module's metal substrate be electrically tied to board GND, or thermally coupled but electrically floating? (Recommend: floating from board GND, on a dielectric, so a single insulation failure does not put 48 V-referenced potential on the heatsink.)

\textbf{8. Does this run own the optics?} PAR-REQ-14 (wide wash beam from 2.5 m, overlapping coverage from 6-8 fixtures) and PAR-REQ-15 (diffusion sufficient that the colours mix before reaching a surface) are stated as requirements on this board, but a PCB run can only choose the emitter and provide mechanical provision for a lens or diffuser. (a) This run selects the emitter and provides lens-holder mounting provision; the diffuser/lens is an enclosure item chosen later. (b) This run owns the full optical chain including the lens/diffuser part numbers and beam-angle verification. \textbf{RECOMMEND (a)} - a PCB pipeline cannot verify a beam angle, and PAR-REQ-14 / 15 should be carried as enclosure requirements in the design doc so they are not silently lost.

\textbf{9. What does "5-10 \% of full output" mean in PAR-REQ-01?} This changes the required driver PWM bandwidth by roughly 40x and it is the single most likely reason this board fails its own review gate 2. At 13-bit / 9.766 kHz the period is 102.4 us and one LSB is \textbf{12.5 ns} - no constant-current driver on the market resolves that, so the real dimming floor is the driver's PWM settling time, not the PWM word length. (a) \textbf{5-10 \% of PWM duty}: on-time 5.1-10.2 us. Most switching CC drivers manage this at 10 kHz with some amplitude error. (b) \textbf{5-10 \% of perceived brightness} (i.e. after gamma \textasciitilde{}2.2): duty is 0.137-0.63 \%, on-time \textbf{141-646 ns}, which is at or below the settling time of essentially every CC driver and is where "visible stair-stepping during slow fades" is actually born. \texttt{CORRECTED 2026-07-28 (orchestrator, verified):} this line previously read "on-time 1.4-6.1 us", a 10x error. The duties were right; the on-times were not. 0.05\textasciicircum{}2.2 x 102.4 us = 140.6 ns and 0.10\textasciicircum{}2.2 x 102.4 us = 646.1 ns. The correction is load-bearing: it moves (b) from "several drivers manage it" to "no surveyed CC driver manages it with PWM alone" - the best part found in P1 (TPS92515HV) specs a 0.2 \% minimum duty at 10 kHz, i.e. a 200 ns floor, which is above the 141 ns that 5 \% perceived requires. See \texttt{research/led-driver.md} and \texttt{research/spec-dimming.md}. \textbf{RECOMMEND: assume (b)}, because that is what a lighting person means by "dim to 5 \%", and it is the harder requirement - designing for (b) satisfies (a) automatically. Please confirm, because it materially narrows the driver shortlist and is the requirement PAR-REQ-09 is really about.

\textbf{10. Calibration path (PAR-REQ-17) - and who does the measuring?} The premise that this is EEPROM \emph{versus} the resistor divider is \textbf{not what the ICD says}. ICD s2 and s3.3: the carrier fits the top leg of the \texttt{ID\_ADC} divider and the daughter fits the bottom leg \textbf{and} an I2C EEPROM may ride the shared I2C bus; a second dedicated ID pin was deleted precisely because the EEPROM does not need one. So the two mechanisms coexist: {*}{*}the divider says what board type this is, the EEPROM stores per-unit calibration.{*}{*} A 24C02/24C32-class part is a few cents and one address allocation (this board owns the whole I2C address space - the carrier reserves nothing - so EEPROM and any digital temp sensor must be allocated non-colliding addresses). \textbf{RECOMMEND: fit both.} The real open question is downstream: PAR-REQ-06 requires measured per-channel scaling, so {*}{*}what instrument measures 6-8 fixtures?{*}{*} (colorimeter / spectrometer / phone-camera comparison / none). If the answer is "nothing", the EEPROM ships empty and colour matching falls back entirely on binning (question 11), which changes the sourcing strategy.

\textbf{11. LED binning versus the assembly route.} PAR-REQ-16 requires binning specified in the BOM for all four channels; "unbinned parts across 8 fixtures will not satisfy PAR-REQ-06". But JLC's catalogue does not generally let you specify a chromaticity/flux bin, so binned high-power emitters usually mean hand-sourcing from a distributor and hand-placing them. (a) Hand-source binned emitters and place them by hand (or as a separate LED sub-assembly), everything else JLC PCBA. (b) Accept whatever bin JLC ships and rely entirely on per-fixture calibration (question 10) to match the fixtures. (c) Both: loose binning plus calibration. \textbf{RECOMMEND (c)} - buy the tightest bin available from a distributor for one batch of 6-8 fixtures (they will then be same-reel and closely matched anyway) \textbf{and} fit the calibration EEPROM. Note (b) alone is only viable if question 10 has a real measurement instrument behind it.

\textbf{12. Confirm the board outline - MECH-02, needed before P5.} The ICD proposes the common footprint: {*}{*}100.0 x 80.0 mm, 3.0 mm corner radius, 1.6 mm thick, 4x M3 at 5 mm inset plus a 5th M3 at (46, 74), and the mandatory 30 x 26 mm notch at (6, 0)-(36, 26)\textbf{. ai-ee has }no outline-shrink step{*}{*} - whatever is passed to \texttt{board\_init} at P5 is permanent, so this must be confirmed, not inherited silently. \textbf{RECOMMEND: confirm 100.0 x 80.0 mm as-is.} The carrier and daughters share one enclosure and mate through the connector, so the outlines are not independent, and after the four exclusion zones and the notch this board has roughly 60 x 74 mm of genuinely usable area for drivers plus whatever the LED scheme needs. Please confirm explicitly, including whether the enclosure has any dimension this footprint must fit inside that has not been stated.

\textbf{13. Assembly method and process.} Not stated anywhere in the LUMINA briefs. (a) JLCPCB fabrication + PCBA, JLC Basic parts preferred where they exist (note: the ICD connectors are Extended, unavoidably - there is no Basic board-to-board part). (b) Hand assembly. Two sub-questions that need answering with it: {*}{*}are the two 2.54 mm sockets hand-soldered post-PCBA{*}{*} (they are THT and must be reverse-mounted on the bottom side, which is not a normal JLC process), and {*}{*}are high-power emitters hand-placed{*}{*} (question 11)? {*}{*}RECOMMEND: JLC PCBA for everything it can do, with the two sockets and the emitters hand-soldered afterward.{*}{*} 6-8 boards is a quantity where two hand operations per board is a couple of hours total, and it removes both awkward processes from the critical path.

\textbf{14. Build quantity and unit cost target for THIS board.} The brief says 6-8 pars of 8-12 fixtures; no per-board cost target is stated (the \$30/board figure in ICD s1.1 belongs to the carrier). The \$500-1000 system budget also predates D-01's at-capable power stage, and the upgrade path implies a \textbf{PoE+} switch, which is materially more expensive than the PoE switch that budget assumed. {*}{*}RECOMMEND: confirm 8 boards built (6-8 deployed plus spares) and set an explicit BOM target for this board.{*}{*} A defensible starting number is \textbf{\$25-35/board} excluding the LED module and heatsink, with the emitter and heatsink budgeted separately because question 2 can double them.

\textbf{15. `ID\_ADC` code allocation.} ICD s3.3: "Board-type codes are allocated by the carrier owner, not chosen by daughters." This board needs its allocated bottom-leg resistor value (and the resulting ADC code) before P4 can place the part. \textbf{RECOMMEND: request the LUM-PAR-A code from the carrier owner now}; if none exists yet, ask for one to be allocated rather than picking a value here - picking one is exactly the silent divergence ICD-01's preamble forbids.

\textbf{16. Sequencing against ICD s7.2 (the one un-frozen section).} The J3/J4 coordinates in s7.2 are provisional until the carrier owner confirms them after its own P6 and re-issues the ICD, and daughters are explicitly blocked on that. Everything else in the ICD is frozen. (a) Wait for the re-issued ICD before starting this run. (b) {*}{*}Proceed through P1-P4 (architecture and schematic, which do not depend on mm coordinates) and hold at P5{*}{*}, where the outline and connector placement are bound permanently. \textbf{RECOMMEND (b)}, with the s7.2 coordinates treated as unconfirmed everywhere they appear, and an explicit re-check against the re-issued ICD before \texttt{board\_init} is run. Please confirm this is acceptable, and if the carrier's P6 has already completed, point this run at the re-issued ICD instead of the DRAFT-A copy in \texttt{brief/}.

\section{Architecture}
\subsection*{Decisions of record (state.json)}
{\small
\begin{longtable}{lp{5.2cm}p{7.2cm}}
\toprule
\textbf{Phase} & \textbf{What} & \textbf{Why} \\
\midrule
\endhead
P0 & D-01/D-02/D-03 inherited as closed; par power budget 8.6-9.3 W af / 18.7-20.0 W at & 05-lumina-closed-decisions.md is binding and supersedes the \textasciitilde{}8.5 W figure in briefs 00/03 (derived from Si3402-B, a part this system does not use) \\
P0 & ICD-01 (brief/06-connector-icd.md) treated as a frozen hard input; no local variant permitted & ICD s10 change process: a daughter that cannot meet it raises a blocking issue against LUM-CAR-A, it does not implement a variant \\
P0 & Common footprint adopted from ICD s7.1: 100.0 x 80.0 mm, R3.0 corners, 5x M3 (4 corners at 5 mm inset + H5 at (46,74)), 1.6 mm FR4, with a mandatory 30x26 mm RJ45 notch at (6,0)-(36,26) & H1-Q4 closed the stack at 11.0 mm; the carrier's board-edge RJ45 is \textasciitilde{}15 mm tall so the notch is a hard mechanical stop and the anti-180-degree interlock (ICD s7.4) \\
P2 & Re-baselined onto ICD rev A2 (copied from boards/lumina-carrier/architecture/connector-icd.md) & Two changes bind this board: (1) board-wide HV clearance raised 0.60 -\textgreater{} 0.635 mm per TPS2378 vendor layout guidance, daughters inherit it; (2) new s8.4 - the carrier's TPS16630 load switch has a mandatory 10k SHDN pull-down, so a daughter may assume an unprogrammed/crashed/reset/brownout carrier presents 0 V at J3, not 48 V. Rail limits s6.2 and mechanical s7.1/s7.2 are UNCHANGED. \\
P2 & LED stage stays on +12V (baseline); the P1 '48V is 16x better for dimming slew' argument is WITHDRAWN as an artifact of holding L fixed & Verified algebraically by the orchestrator: with L sized for the same ripple fraction k, t\_rise+t\_fall = 1/(f{*}k), independent of Vin. At 700 kHz / k=0.30 both rails give 4.762 us (48V only redistributes 3.81/0.95 us between rise and fall). That is 33.8x the 141 ns PAR-REQ-01 budget, so NO rail or inductor choice reaches the requirement by gating the converter - it needs shunt-FET dimming either way. The 48V branch therefore buys nothing on fidelity and costs 0.635 mm creepage, 100 V caps, bleed path and daughter-owned inrush. \\
\bottomrule
\end{longtable}
}
\subsection*{Sheet plan}
\subsection*{lumina-par (LUM-PAR-A) - hierarchical sheet plan}
Five sheets under a root that contains nothing but the sheet symbols and the inter-sheet wiring. Interface nets become hierarchical pins at P4 and placement groups at P6.

\textbf{Net-naming contract (binding on P4).} KiCad names a net from a hierarchical label alone as \texttt{/\textless{}sheet\textgreater{}/\textless{}LABEL\textgreater{}}. To produce the canonical names below, {*}{*}P4 must place a root-sheet local label on every inter-sheet wire{*}{*}, spelled exactly as in s2. \texttt{constraints.json} uses these names and P5-P8 will silently no-op on any net whose name does not match. Power nets are bare global symbols (\texttt{+3V3}, \texttt{+12V}, \texttt{+48V\_SW}, \texttt{GND}); every other inter-sheet net carries a leading \texttt{/}.

---

\subsubsection*{1. Sheets, refdes ranges and contents}
Refdes are unique across sheets by construction: each sheet owns a hundreds block for every prefix, and each sheet has its own \texttt{\#PWR} base.

\begin{quote}\small\ttfamily\raggedright
\textbar{} Sheet \textbar{} Blocks \textbar{} Refdes block \textbar{} `\#PWR` base \textbar{}\\
\textbar{}---\textbar{}---\textbar{}---\textbar{}---\textbar{}\\
\textbar{} `power` \textbar{} B1 - rail entry, bulk, branch-B front end (DNP), test points, H5 \textbar{} {*}{*}100-199{*}{*} \textbar{} {*}{*}100{*}{*} \textbar{}\\
\textbar{} `control` \textbar{} B2 - ENABLE gating and fail-safe logic; B5 - ID and calibration \textbar{} {*}{*}200-299{*}{*} \textbar{} {*}{*}200{*}{*} \textbar{}\\
\textbar{} `drivers` \textbar{} B3 - the four constant-current channels and their shunt FETs \textbar{} {*}{*}300-399{*}{*} \textbar{} {*}{*}300{*}{*} \textbar{}\\
\textbar{} `thermal` \textbar{} B4 - NTC dividers, window comparator, FAULT driver \textbar{} {*}{*}400-499{*}{*} \textbar{} {*}{*}400{*}{*} \textbar{}\\
\textbar{} `led\_if` \textbar{} B6 - LED harness header, TVS clamps, thermocouple pads \textbar{} {*}{*}500-599{*}{*} \textbar{} {*}{*}500{*}{*} \textbar{}
\end{quote}
\textbf{Exception: J3, J4 and H5 keep the refdes the ICD assigns them} (ICD s7.2 / s7.5 name them J3, J4 and H5, and the carrier's placement tables refer to them by those names). They are exempt from the hundreds blocks. J5 is this board's own harness header and follows the same connector sequence.

\paragraph*{1.1 \texttt{power} (100-199, \texttt{\#PWR} 100)}
\begin{quote}\small\ttfamily\raggedright
\textbar{} Ref \textbar{} Part \textbar{} Note \textbar{}\\
\textbar{}---\textbar{}---\textbar{}---\textbar{}\\
\textbar{} {*}{*}J3{*}{*} \textbar{} CONNFLY DS1023-2{*}7SF11, 2x7 socket \textbar{} {*}{*}reverse-mounted, bottom side, facing down{*}{*} (ICD s7.3). Pin 1 silkscreened with a triangle \textbar{}\\
\textbar{} {*}{*}H5{*}{*} \textbar{} MountingHole\_3.2mm\_M3 \textbar{} (46, 74), ICD s7.5. Not a `board\_init` hole - added here so it carries a refdes \textbar{}\\
\textbar{} C101, C102 \textbar{} 22 uF / 25 V X7R \textbar{} `+12V` bulk \textbar{}\\
\textbar{} C103, C104 \textbar{} 10 uF / 16 V X7R \textbar{} `+3V3` bulk \textbar{}\\
\textbar{} C105 \textbar{} 100 nF \textbar{} `+3V3` HF \textbar{}\\
\textbar{} TP101-TP103 \textbar{} test pads \textbar{} `+12V`, `+3V3`, `+48V\_SW`. {*}{*}TP103 carries the ICD s9 bench-hazard silkscreen{*}{*} \textbar{}\\
\textbar{} {*}{*}Q101, Q102, R101-R104, C106-C108{*}{*} \textbar{} {*}{*}branch-B front end - DNP on branch A{*}{*} \textbar{} 100 V P-FET hot-swap, Cgd dV/dt network, gate pull-down N-FET from `/EN\_OK`, 100 k {*}{*}0805{*}{*} bleed, 2x 10 uF/100 V bulk. R104 is 0805 now so branch B needs no footprint change (ICD s5.4) \textbar{}
\end{quote}
\paragraph*{1.2 \texttt{control} (200-299, \texttt{\#PWR} 200)}
\begin{quote}\small\ttfamily\raggedright
\textbar{} Ref \textbar{} Part \textbar{} Note \textbar{}\\
\textbar{}---\textbar{}---\textbar{}---\textbar{}\\
\textbar{} {*}{*}J4{*}{*} \textbar{} CONNFLY DS1023-2{*}12SF11, 2x12 socket \textbar{} reverse-mounted, bottom side, facing down \textbar{}\\
\textbar{} U201 \textbar{} SN74LVC1G08 (SOT-23-5) \textbar{} `/EN\_OK = ENABLE AND FAULT` \textbar{}\\
\textbar{} U202 \textbar{} 74LVC00A quad NAND (SO-14) \textbar{} `/SHUNTn = NOT(PWMn AND /EN\_OK)` - one package covers all four channels \textbar{}\\
\textbar{} U203 \textbar{} 24C32-class I2C EEPROM \textbar{} {*}{*}address 0x50{*}{*}, WP to GND. {*}{*}No I2C pull-ups on this board{*}{*} \textbar{}\\
\textbar{} R201 \textbar{} 100 k \textbar{} {*}{*}ENABLE pull-down at the connector end, before any series element{*}{*} (ICD s8.2) \textbar{}\\
\textbar{} R202-R205 \textbar{} 4x 100 k \textbar{} PWM0-3 pull-downs at J4 - an undriven carrier must not float a gate input high \textbar{}\\
\textbar{} R206 \textbar{} {*}{*}VALUE TBD{*}{*} \textbar{} `ID\_ADC` bottom leg. {*}{*}Allocated by the carrier owner, not chosen here{*}{*} (ICD s3.3). P4 blocker - `decisions.md` OPEN-2 \textbar{}\\
\textbar{} R207 \textbar{} 100 k \textbar{} local `FAULT` pull-up to `+3V3`, so the node is defined when unmated \textbar{}\\
\textbar{} R208, R209 \textbar{} 0 R / 10 k \textbar{} EEPROM A0-A2 and WP strapping \textbar{}\\
\textbar{} R214-R217 \textbar{} 4x 0 R \textbar{} `/EN\_OK -\textgreater{} /DRV\_ENn` links, {*}{*}fitted by default{*}{*} \textbar{}\\
\textbar{} {*}{*}U204, R210-R213, C210-C213, D201-D204{*}{*} \textbar{} {*}{*}converter-idle one-shot - DNP option{*}{*} \textbar{} 74LVC14 hex Schmitt + retriggerable RC + diode per channel. Populate to save \textasciitilde{}0.8 W at saturated colours (`power\_tree.md` s2); if populated, remove R214-R217 \textbar{}\\
\textbar{} C201-C205 \textbar{} decoupling + `ID\_ADC` 100 nF \textbar{} \textbar{}
\end{quote}
\paragraph*{1.3 \texttt{drivers} (300-399, \texttt{\#PWR} 300)}
Four identical channels; channel \texttt{n} is based at \texttt{300 + 20n}. Channel 0 shown; channels 1-3 are U321/U341/U361 and so on.

\begin{quote}\small\ttfamily\raggedright
\textbar{} Ref \textbar{} Part \textbar{} Note \textbar{}\\
\textbar{}---\textbar{}---\textbar{}---\textbar{}\\
\textbar{} U301 \textbar{} TPS92515HVDGQR class (MSOP-10-EP) \textbar{} 5.5-65 V, 2 A internal FET, shunt-PWM capable. Exposed pad on a \textgreater{}= 9-via array \textbar{}\\
\textbar{} L301 \textbar{} 47 uH shielded \textbar{} sized for {*}{*}ripple, not slew{*}{*} - 90 mA (30 \%) at \textasciitilde{}700 kHz \textbar{}\\
\textbar{} D301 \textbar{} {*}{*}60 V{*}{*} Schottky, 1 A \textbar{} branch A. Low Vf matters: it dominates the 0.20 W idle loss. Branch B needs a 100 V part \textbar{}\\
\textbar{} R301 \textbar{} sense resistor \textbar{} sets 300 mA (af, branch A). {*}{*}The at upgrade is this value{*}{*} \textbar{}\\
\textbar{} Q301 \textbar{} 60 V logic-level N-FET \textbar{} shunt across the string. {*}{*}60 V now, not 30 V{*}{*}, so the branch-B / `at` 8S string (27.2 V) needs no change \textbar{}\\
\textbar{} R302 \textbar{} 150 R \textbar{} shunt gate series - bounds the 74LVC00 output to 22 mA, \textasciitilde{}30 ns edges \textbar{}\\
\textbar{} R303 \textbar{} 100 k \textbar{} shunt gate pull-down - an unpowered gate IC must not leave a FET half-on \textbar{}\\
\textbar{} C301 \textbar{} 4.7 uF / 25 V \textbar{} VIN local (100 V on branch B) \textbar{}\\
\textbar{} C302 \textbar{} 100 nF \textbar{} VIN HF, at the pin \textbar{}\\
\textbar{} C303, R304 \textbar{} {*}{*}DNP{*}{*} \textbar{} SW-node RC snubber footprint \textbar{}
\end{quote}
\textbf{There is no output capacitor across any LED string} - a shunt FET dumps it every PWM cycle (\texttt{blocks.md} B3 rule 1). Do not let a reviewer add one.

\paragraph*{1.4 \texttt{thermal} (400-499, \texttt{\#PWR} 400)}
\begin{quote}\small\ttfamily\raggedright
\textbar{} Ref \textbar{} Part \textbar{} Note \textbar{}\\
\textbar{}---\textbar{}---\textbar{}---\textbar{}\\
\textbar{} U401 \textbar{} quad open-drain comparator, LMV339 / TLV3704 class \textbar{} rail-to-rail input, \textless{}= 50 uA/channel. {*}{*}Open-drain outputs are mandatory{*}{*} - they wire-OR onto `FAULT`, which must never be driven high \textbar{}\\
\textbar{} RT401 \textbar{} 10 k B3950 NTC, 0603 \textbar{} {*}{*}on the copper of the hottest driver stage, outside the DC-DC hot zone (2,46)-(36,68){*}{*} \textbar{}\\
\textbar{} R401, R403 \textbar{} 2x 10 k \textbar{} NTC divider top legs (board, module) \textbar{}\\
\textbar{} R402, R404 \textbar{} \textless{}= 1 k \textbar{} ADC series - keeps total source impedance \textless{}= 6 kohm against the ICD's 10 kohm \textbar{}\\
\textbar{} R405-R412 \textbar{} reference ladder + hysteresis \textbar{} ratiometric off `+3V3`, so thresholds track the carrier's ADC reference \textbar{}\\
\textbar{} C401-C403 \textbar{} decoupling + NTC filter \textbar{} filter caps sized so R402/R404 x C stays a slow-signal filter, not a pulse filter \textbar{}
\end{quote}
Four comparator channels: {*}{*}emitter hot (90 C, release 75 C) / emitter open / emitter short / board hot (110 C, release 95 C){*}{*}. The window (open + short) converts the most likely fault in an off-board module - a broken harness wire, which reads as \emph{cold} in either divider orientation - from fail-dangerous to fail-safe (D-T17, E-10).

\paragraph*{1.5 \texttt{led\_if} (500-599, \texttt{\#PWR} 500)}
\begin{quote}\small\ttfamily\raggedright
\textbar{} Ref \textbar{} Part \textbar{} Note \textbar{}\\
\textbar{}---\textbar{}---\textbar{}---\textbar{}\\
\textbar{} J5 \textbar{} 10-way latched wire-to-board, 2.0 mm (JST PH class) \textbar{} 4 anode + 4 per-channel GND return + `/NTC\_LED` + dedicated NTC sense return \textbar{}\\
\textbar{} D501-D504 \textbar{} TVS, SOD-123 \textbar{} {*}{*}15 V standoff on branch A, 33 V on branch B{*}{*} - same footprint, a BOM value change \textbar{}\\
\textbar{} TP501, TP502 \textbar{} bare-copper thermocouple pads \textbar{} bring-up thermal verification (L-13). ICD s9 bench-hazard silkscreen \textbar{}
\end{quote}
\textbf{No 48 V net reaches this sheet, ever} (D-T13). That is what keeps the 0.635 mm clearance rule, the 100 V capacitor rule and the 0805-minimum resistor rule off the harness and the module.

---

\subsubsection*{2. Interface nets - canonical names}
These are the names the schematic \textbf{must} produce. \texttt{constraints.json} and every P5-P8 check key off them.

\begin{quote}\small\ttfamily\raggedright
\textbar{} Net \textbar{} Type \textbar{} From -\textgreater{} to \textbar{} Note \textbar{}\\
\textbar{}---\textbar{}---\textbar{}---\textbar{}---\textbar{}\\
\textbar{} `+48V\_SW` \textbar{} power \textbar{} `power` (J3 1/3/5) \textbar{} {*}{*}landed, not tapped{*}{*} on branch A. Declared in `voltages` at 57 V so `check\_creepage` enforces the clearance \textbar{}\\
\textbar{} `+12V` \textbar{} power \textbar{} `power` -\textgreater{} `drivers` \textbar{} 0.718 A max \textbar{}\\
\textbar{} `+3V3` \textbar{} power \textbar{} `power` -\textgreater{} `control`, `thermal` \textbar{} \textless{}= 5 mA \textbar{}\\
\textbar{} `GND` \textbar{} power \textbar{} everywhere \textbar{} In1 plane \textbar{}\\
\textbar{} `/PWM0` .. `/PWM3` \textbar{} in \textbar{} `control` (J4) internal \textbar{} 3.3 V CMOS. {*}{*}No RC filter; if any network, tau \textless{}= 14 ns{*}{*} \textbar{}\\
\textbar{} `/ENABLE` \textbar{} in \textbar{} `control` (J4) internal \textbar{} 100 k pull-down here \textbar{}\\
\textbar{} `/FAULT` \textbar{} bidir, open drain \textbar{} `thermal` -\textgreater{} `control` -\textgreater{} J4 \textbar{} {*}{*}never driven high{*}{*}; wire-OR'd with the carrier eFuse fault \textbar{}\\
\textbar{} `/EN\_OK` \textbar{} internal \textbar{} `control` -\textgreater{} `power` (branch-B gate) \textbar{} `ENABLE AND FAULT` \textbar{}\\
\textbar{} `/DRV\_EN0` .. `/DRV\_EN3` \textbar{} out \textbar{} `control` -\textgreater{} `drivers` \textbar{} driver PWM/EN pins. Tied to `/EN\_OK` through 0 R by default; driven by the one-shot if populated \textbar{}\\
\textbar{} `/SHUNT0` .. `/SHUNT3` \textbar{} out \textbar{} `control` -\textgreater{} `drivers` \textbar{} shunt-FET gates. `NOT(PWMn AND /EN\_OK)` \textbar{}\\
\textbar{} `/LED0\_A` .. `/LED3\_A` \textbar{} out \textbar{} `drivers` -\textgreater{} `led\_if` \textbar{} channel anodes, 0.30 A each \textbar{}\\
\textbar{} `/NTC\_LED` \textbar{} in \textbar{} `led\_if` -\textgreater{} `thermal` \textbar{} module NTC divider node \textbar{}\\
\textbar{} `/NTC\_BRD` \textbar{} internal \textbar{} `thermal` \textbar{} board NTC divider node \textbar{}\\
\textbar{} `/ADC0` \textbar{} out \textbar{} `thermal` -\textgreater{} `control` -\textgreater{} J4 \textbar{} emitter temperature to the carrier. Source impedance \textless{}= 6 kohm \textbar{}\\
\textbar{} `/ADC1` \textbar{} out \textbar{} `thermal` -\textgreater{} `control` -\textgreater{} J4 \textbar{} board temperature to the carrier \textbar{}\\
\textbar{} `/ID\_ADC` \textbar{} out \textbar{} `control` -\textgreater{} J4 \textbar{} divider bottom leg; value allocated by the carrier owner \textbar{}\\
\textbar{} `/I2C\_SCL`, `/I2C\_SDA` \textbar{} bidir \textbar{} `control` (J4) -\textgreater{} U203 \textbar{} {*}{*}no daughter-side pull-ups{*}{*} \textbar{}
\end{quote}
\textbf{Unused ICD signals}, landed on J4 and left unconnected on this board: \texttt{PWM4..PWM7} (no fifth channel on rev A), \texttt{DSPI\_SCK/MOSI/MISO/CSn} (a 4-channel PWM design needs no serial device beyond the EEPROM, and \texttt{decisions.md} D2 rejects a local PWM generator). Leave them as no-connects with an explicit \texttt{\textasciitilde{}} ERC directive so P4's ERC does not flag them and so a reviewer can see the decision.

\textbf{I2C address map} (this board owns the whole space; the carrier reserves nothing): {*}{*}0x50 = 24C32 EEPROM; 0x51-0x57 reserved for its own address pins; nothing else on the bus.{*}{*} The comparator and both NTCs are analogue.

---

\subsubsection*{3. Placement groups for P6}
Declared in \texttt{constraints.json} under \texttt{placement.groups}. Each buck channel is one tight cluster; the loop that matters is \texttt{+12V} -\textgreater{} C301 -\textgreater{} U301 -\textgreater{} L301 -\textgreater{} D301.

\begin{quote}\small\ttfamily\raggedright
\textbar{} Group \textbar{} Anchor \textbar{} Members \textbar{}\\
\textbar{}---\textbar{}---\textbar{}---\textbar{}\\
\textbar{} `ch0` \textbar{} U301 \textbar{} L301, D301, R301, R302, R303, C301, C302, Q301 \textbar{}\\
\textbar{} `ch1` \textbar{} U321 \textbar{} L321, D321, R321, R322, R323, C321, C322, Q321 \textbar{}\\
\textbar{} `ch2` \textbar{} U341 \textbar{} L341, D341, R341, R342, R343, C341, C342, Q341 \textbar{}\\
\textbar{} `ch3` \textbar{} U361 \textbar{} L361, D361, R361, R362, R363, C361, C362, Q361 \textbar{}\\
\textbar{} `ntc\_brd` \textbar{} U301 \textbar{} RT401 - the board NTC must sit {*}{*}on{*}{*} the hottest driver stage (T-3) \textbar{}\\
\textbar{} `gate` \textbar{} U202 \textbar{} R202-R205, C202 - keep the NAND next to J4's PWM pins to keep the jitter-sensitive run short \textbar{}\\
\textbar{} `analog` \textbar{} U401 \textbar{} R401-R412, C401-C403 - {*}{*}separated from every inductor by \textgreater{}= 10 mm{*}{*} \textbar{}
\end{quote}
\texttt{placement.fixed}: \textbf{H5}, positioned by an explicit \texttt{place\_edit} op at (46, 74) in board-local coordinates (translated - see \texttt{stackup.md} TRAP 2).

\texttt{placement.edges}:

\begin{quote}\small\ttfamily\raggedright
\textbar{} Ref \textbar{} Edge \textbar{} pos \textbar{} Why \textbar{}\\
\textbar{}---\textbar{}---\textbar{}---\textbar{}---\textbar{}\\
\textbar{} J3 \textbar{} bottom \textbar{} 0.24 \textbar{} ICD s7.2 nominal centre x = 24 mm. {*}{*}Hint only{*}{*} \textbar{}\\
\textbar{} J4 \textbar{} bottom \textbar{} 0.72 \textbar{} ICD s7.2 nominal centre x = 72 mm. {*}{*}Hint only{*}{*} \textbar{}\\
\textbar{} J5 \textbar{} {*}{*}top{*}{*} \textbar{} 0.56 \textbar{} x \textasciitilde{} 56 mm - the only clear span of the top edge (the notch occupies x 6-36, the recovery-header keepout x 76-98), and it sits directly above the driver band so the LED anode runs stay short \textbar{}
\end{quote}
\textbf{J3 and J4's entries are hints to the annealer, not final positions.} Both must be placed with explicit \texttt{place\_edit} ops at the re-issued ICD s7.2 coordinates and then \textbf{locked}, before the annealer runs. Both are \textbf{bottom-side, facing down} - set the side at P4 in the footprint or with a \texttt{place\_edit} \texttt{flip} op, and check pin 1 in the mated view, not from the footprint (pin numbering mirrors across the row on a bottom-side part).

---

\subsubsection*{4. Sheet-level P4 notes}
1. \textbf{The DNP option sets must be in the netlist}, not added later: the branch-B front end (\texttt{power} Q101/Q102/R101-R104/C106-C108) and the converter-idle one-shot (\texttt{control} U204/R210-R213/C210-C213/D201-D204). A DNP part still has pads, so its clearance and area are accounted for at P6/P7 and the option stays a populate change rather than a respin. Mark them \texttt{DNP} in the BOM variant field so \texttt{bom\_cpl} at P9 excludes them. 2. \textbf{`FAULT` is open drain.} ERC will want a driving pin somewhere; the comparator outputs plus R207 are it. Do not let a fix loop add a push-pull driver. 3. \textbf{No I2C pull-ups.} Fitting them is an ICD violation, and P4's reflex is to add them. Record it on the schematic. 4. {*}\emph{`PWM4..7` and `DSPI\_}` are deliberate no-connects.{*}{*} 5. \textbf{R206's value is unknown until the carrier owner allocates it.} Place the part with a \texttt{TBD} value field and a schematic note; do not guess a value - that is exactly the silent divergence ICD-01's preamble forbids.

Full architecture narrative (not inlined; see the artifact index): \texttt{architecture/blocks.md}, \texttt{architecture/decisions.md}, \texttt{architecture/power\_tree.md}, \texttt{architecture/stackup.md}.
\section{Schematic}
\emph{Pending --- produced at P4.}
\section{Layout}
\emph{Pending --- produced at P6.}
\section{Verification}
\emph{Pending --- produced at P8.}
\section{DFM and Fabrication}
\emph{Pending --- produced at P9.}
\section{Run Record (Appendix)}
\subsection*{Phase digests}
\paragraph*{P0}
\begin{quote}\small\ttfamily\raggedright
\# P0 Intake digest - lumina-par (LUM-DTR-PAR-A / LUM-PAR-A)\\
~\\
- Board: RGBW par daughter for LUM-CAR-A. 4 (maybe 5) PWM-dimmed constant-current LED\\
  channels, stacked 11.0 mm above the carrier on the frozen ICD-01 connector pair.\\
  Continuous duty; the design problem is dimming quality and fixture-to-fixture colour\\
  consistency, not peak power.\\
- Inherited as CLOSED (05-lumina-closed-decisions.md): D-01 (at-capable stage, af-classified;\\
  budget 8.6-9.3 W af / 18.7-20.0 W at to this daughter - the "\textasciitilde{}8.5 W" and "\textasciitilde{}10 W regulated"\\
  figures in briefs 00/03 are dead), D-02 (48 V + 12 V + 3.3 V on the connector), D-03 (no UV).\\
- ICD-01 is a frozen hard input. All interfaces are ICD-only: J3 2x7 power, J4 2x12 signal.\\
  ICD s9 forbids adding an external connector of any kind. Common footprint 100 x 80 mm,\\
  R3.0, 5x M3, 1.6 mm, plus a mandatory 30 x 26 mm RJ45 notch at (6,0)-(36,26).\\
- Safety flags: \textgreater{}30 V unconditional (48 V present on J3 whether tapped or not - 0.60 mm outer\\
  clearance, 100 V caps, 0805+ resistors across the 48 V domain); whole fixture floats at PoE\\
  potential including LED wiring and any heatsink; \textgreater{}3 A conditional on per-channel LED current\\
  that P1/P2 must publish.\\
- Riskiest unknowns, in order: (1) af-vs-at sizing - +12V tops out at 15.0 W so at-sizing forces\\
  a \textgreater{}=60 V stage on THIS board, the exact thing D-02's 12 V rail existed to avoid; (2) continuous\\
  thermal in a non-conductive enclosure at 56-69 C internal air; (3) PAR-REQ-01's "5-10 \% of full\\
  output" is ambiguous and changes the required driver settling time by \textasciitilde{}40x.\\
- Hard derived requirement recorded: the board must be electrically incapable of exceeding the\\
  carrier's fault ceilings (12 V OCP 2.0 A, 48 V eFuse 1.0 A latch-off) with every PWM stuck at\\
  100 \%, because the PAR-REQ-11 clamp is firmware and firmware hangs.\\
- Pipeline risk logged for P5: `board\_init.py` has no notch/cutout option and `kc.py` has no\\
  outline editor, so the mandatory RJ45 relief needs a direct Edge.Cuts edit after board\_init;\\
  and `--mounting-holes` places holes at margin/2, so the ICD's 5 mm inset needs `--margin 10`.\\
- Artifact: `requirements.md` (9 sections, 16 open questions, all with recommendations).
\end{quote}
\paragraph*{P1}
\begin{quote}\small\ttfamily\raggedright
\# P1 Research digest - lumina-par\\
~\\
Roster: 2x component-scout (LED driver, RGBW emitter), 1x reference-design\\
(continuous thermal + protection chain), 1x interface-spec (PWM dimming /\\
flicker / perceptual resolution). No power-architect: the rail tree is fixed by\\
ICD-01 s6.2 and the only open variable was which rail feeds the LED stage, which\\
became conflict C1 for the architect.\\
~\\
- {*}{*}Thermal is the binding constraint, not the power budget.{*}{*} A sealed\\
  120x100x60 mm non-metallic box is 3.6-4.3 K/W internal-air-to-room (Hoffman\\
  \textasciitilde{}4.6 W/m2K, cross-checked vs Rittal k=5.5). Nothing on the PCB changes that\\
  number. Sealed caps total box heat at \textasciitilde{}5.0 W for 45 C internal air in a 25 C\\
  room; minus the carrier's 2.4 W that leaves \textasciitilde{}2.6 W of daughter heat = \textasciitilde{}3-4 W of\\
  LED electrical power, {*}{*}under half the af envelope{*}{*}.\\
- {*}{*}ICD s7.6's internal-air figures do not hold.{*}{*} Its 69 C (at) is internally\\
  inconsistent with its own 56 C (af): convection is near-linear in dT at this\\
  scale, so \textasciitilde{}2x the heat must give \textasciitilde{}2x the rise (31 K -\textgreater{} \textasciitilde{}60 K -\textgreater{} \textasciitilde{}85 C), and the\\
  independent calculation gives 89-115 C. Optimistic by 20-46 K. Raised as a\\
  blocking issue against LUM-CAR-A under ICD s10 (OPEN-1).\\
- {*}{*}PAR-REQ-01 at 5 \% perceived fails at 13-bit by 8.9x{*}{*} and needs \textasciitilde{}17 bits;\\
  ESP32-S3 LEDC caps at 14. Two independent walls: resolution (17 bits) and\\
  driver settling (T \textless{}= 141 ns). New fact, verified from ESP-IDF source: LEDC does\\
  NOT rigidly couple frequency to resolution - "80 MHz / 2\textasciicircum{}n" is a maximum, so\\
  14-bit is available well below 4.88 kHz. Widens the trade space.\\
- {*}{*}Flicker is over-specified.{*}{*} IEEE 1789-2015 RP2 removes all modulation-depth\\
  restriction above 3 kHz; 9.766 kHz sits 3.3x above it. The 10 kHz-over-1.2 kHz\\
  camera call is confirmed (8.1x lower band contrast) but bounded: clean to\\
  120 fps, marginal at 240 fps, 960 fps unreachable at any usable LEDC setting.\\
- {*}{*}No CC LED driver has an enable pin separate from its PWM pin.{*}{*} On all five\\
  surviving candidates the dimming pin IS the enable, so ENABLE must be gated with\\
  each PWM line externally. The brief's "driver with an EN pin strongly preferred"\\
  is not achievable in this performance class.\\
- {*}{*}Emitter is single-source.{*}{*} The XINGLIGHT 4-in-1 is the only in-stock RGBW\\
  power emitter above 50 mA on the whole LCSC catalogue (318-part sweep; every\\
  Cree XLamp colour line at 0 stock), with no published Rth and no published Tj\\
  max. Fallback is RGB 3-in-1 + separate white at the cost of PAR-REQ-15.\\
- Orchestrator correction applied to `requirements.md`: s9 Q9(b) on-times were 10x\\
  high (1.4-6.1 us -\textgreater{} 141-646 ns), verified by independent recompute. Load-bearing:\\
  it moved the strict reading of PAR-REQ-01 from "several drivers manage it" to\\
  "no surveyed CC driver manages it by gating the converter".
\end{quote}
\paragraph*{P2}
\begin{quote}\small\ttfamily\raggedright
\# P2 Architecture digest - lumina-par\\
~\\
Artifact: `architecture/` (blocks.md, power\_tree.md, stackup.md, sheets.md,\\
constraints.json, decisions.md). This is the H1 checkpoint package.\\
~\\
- {*}{*}6 blocks{*}{*}: rail entry (48 V landed on J3 but NOT tapped), ENABLE gating,\\
  4x TPS92515HV-class CC buck with a shunt FET per channel, window-comparator\\
  over-temperature protection, ID divider + I2C EEPROM, 10-way LED harness.\\
  Emitters off-board on an aluminium MCPCB.\\
- {*}{*}Stackup JLC04161H-3313, 4 layers{*}{*}, 100.0 x 80.0 mm, 1.6 mm. Drivers: switch-node\\
  containment 11 mm above a live 2.4 GHz antenna, a continuous GND reference under\\
  jitter-sensitive PWM, B.Cu consumed by two reverse-mounted sockets, \textasciitilde{}120 parts in\\
  \textasciitilde{}45 cm2. No controlled-impedance net exists on this board.\\
- {*}{*}Design point 150 mA/die, 4 packages, 2S2P per channel.{*}{*} Full white (the\\
  stuck-PWM hardware backstop) = 0.718 A on +12V: 4 \% under the sustained ceiling\\
  and 2.8x under OCP. The board is electrically incapable of exceeding its budget.\\
  The \textgreater{}3 A conditional flag is CLOSED - max per-channel current is 0.30 A.\\
- {*}{*}C1 arbitrated to +12V{*}{*}, and the P1 argument for 48 V was withdrawn rather than\\
  overruled. The "48 V gives 16x faster slew" claim held the inductor fixed across a\\
  4x change in (Vin - Vs). With L sized for the same ripple fraction, t\_r + t\_f =\\
  1/(f{*}k) - the rail voltage cancels out of the algebra. Verified independently by\\
  the orchestrator: 4.762 us on BOTH rails at 700 kHz / k=0.30 (48 V only\\
  redistributes 3.81/0.95 us between rise and fall). That is 33.8x the 141 ns\\
  budget, so no rail or inductor choice reaches PAR-REQ-01 by gating the converter.\\
- {*}{*}C2 resolved: shunt-FET dimming closes the timing wall (T \textasciitilde{} 50-100 ns vs 141 ns,\\
  1.4-2.8x margin); the resolution wall is NOT closable on this board{*}{*} and needs\\
  carrier firmware dithering (CR-5). spec-dimming's "local \textgreater{}=16-bit PWM IC" class\\
  LOSES to led-driver's six-part survey of that exact class - no part delivers\\
  \textgreater{}=150 mA/ch AND \textgreater{}=10 kHz AND useful resolution AND independent channels.\\
- {*}{*}Cost \textasciitilde{}\$18-23/board{*}{*} delivered (BOM \textasciitilde{}\$12.9, driver silicon 52 \%) against a\\
  \$25-35 target; module + heatsink adds \$8-14/fixture.\\
- {*}{*}Sealed enclosure does not close{*}{*} with the module inside (9.15 W -\textgreater{} 62 C air\\
  against a 45 C criterion). It closes vented, or sealed with LED heat conducted\\
  through the wall (3.21 W -\textgreater{} 38 C). Seven ENC-x acceptance criteria written for\\
  whoever owns the enclosure, since it is not this board's deliverable.\\
- Riskiest decision: {*}{*}emitters off-board on MCPCB (D3){*}{*} - the thermal argument\\
  rests on an ASSUMED 12 K/W package Rth, because the vendor publishes neither Rth\\
  nor Tj max, and it is the only in-stock RGBW power emitter on LCSC. At 20 K/W even\\
  MCPCB is marginal at af.\\
- Two P5 traps recorded: fixed-outline board\_init derives its origin from the packed\\
  component bbox (every rect in constraints.json needs a P5 translation), and the\\
  antenna column has no automated checker on F.Cu/B.Cu.
\end{quote}
\subsection*{Gate attempt history}
\emph{(no entries)}
\subsection*{Issues}
\emph{(no entries)}
\subsection*{Human checkpoints}
{\small
\begin{longtable}{lllp{2.6cm}p{6.4cm}}
\toprule
\textbf{Id} & \textbf{Phase} & \textbf{Status} & \textbf{When} & \textbf{Note} \\
\midrule
\endhead
1 & P2 & not reached &  &  \\
2 & P4 & not reached &  &  \\
3 & P6 & not reached &  &  \\
4 & P8 & not reached &  &  \\
5 & P10 & not reached &  &  \\
\bottomrule
\end{longtable}
}
\section{Artifact Index}
{\small
\begin{longtable}{p{9cm}rp{4cm}}
\toprule
\textbf{File} & \textbf{Size} & \textbf{Note} \\
\midrule
\endhead
\texttt{state.json} & 5,826 B &  \\
\texttt{requirements.md} & 39,521 B &  \\
\texttt{architecture/blocks.md} & 28,172 B &  \\
\texttt{architecture/constraints.json} & 7,703 B &  \\
\texttt{architecture/decisions.md} & 34,689 B &  \\
\texttt{architecture/power\_tree.md} & 17,385 B &  \\
\texttt{architecture/sheets.md} & 12,214 B &  \\
\texttt{architecture/stackup.md} & 24,580 B &  \\
\bottomrule
\end{longtable}
}
\end{document}
